% Based on: multi-query kinodynamic planning paper.

\chapter{Multi-Query Kinodynamic Planning via Lazy Edge Bundles}
\label{ch:multiquery}

% TODO: Open with the bottleneck that motivates a multi-query approach to
% kinodynamic planning.

\section{The Bottleneck}
\label{sec:multiquery:bottleneck}

% TODO: Kinodynamic planning is slow in high dimensions; single-query
% approaches rebuild the search tree per query; simulation and collision
% checking are expensive (pre-VAMP, pre-cuRobo).

\section{Method: Edge Bundles over Discretized State and Control Spaces}
\label{sec:multiquery:method}

\subsection{Edges as Sampled Velocity Sequences}
\label{sec:multiquery:edges}

% TODO: Define an edge as (start config, goal config, sampled velocity
% sequence, simulation duration).

\subsection{Disconnected Bundles, Not Connected Roadmaps}
\label{sec:multiquery:bundles}

% TODO: Contrast the disconnected bundle structure with PRM-style
% connected graphs.

\subsection{Probabilistic and Asymptotic Completeness}
\label{sec:multiquery:completeness}

% TODO: Establish completeness via random duration sampling.

\section{Lazy Evaluation via Neighborhood Structure}
\label{sec:multiquery:lazy}

\subsection{Lazy Propagation Probability}
\label{sec:multiquery:propagation}

% TODO: Derive the lazy propagation probability from theta-neighborhood
% edge counts.

\subsection{Lazy Collision Probability}
\label{sec:multiquery:collision}

% TODO: Analogous formulation for lazy collision checking.

\section{The Planner}
\label{sec:multiquery:planner}

% TODO: RRT-style or F-score node selection, with a lazy-then-actual
% propagation strategy.

\section{Results and Trade-Offs}
\label{sec:multiquery:results}

\subsection{Amortization with Known Dynamics and Collision Geometry}
\label{sec:multiquery:amortization}

% TODO: Quantify how far amortization goes when dynamics and collision
% geometry are known up front.

\subsection{Pre-Processing Cost and Re-Bundling}
\label{sec:multiquery:preprocess}

% TODO: Pre-processing cost, and the re-bundling required when dynamics
% change.

\subsection{Path Quality vs.\ Storage and Lookup Costs}
\label{sec:multiquery:quality}

% TODO: How path quality scales against storage and lookup cost as
% sampling density grows.

\section{The Opening This Leaves}
\label{sec:multiquery:opening}

% TODO: Pose the question that bridges to Chapter~\ref{ch:generation}:
% can pre-computation scale further via learning?
